\documentclass[UTF8]{ctexart}
\usepackage[a4paper,margin=2.5cm]{geometry}
\usepackage{booktabs}
\usepackage{hyperref}
\usepackage{listings}

\title{ASC 队员选拔最终报告}
\author{姓名}
\date{}

\begin{document}
\maketitle

\section{基本信息}

\begin{tabular}{ll}
\toprule
项目 & 内容 \\
\midrule
姓名 &  \\
年级专业 &  \\
基础题 & HPL / HPCG \\
大题 &  \\
基础题仓库链接 &  \\
大题仓库链接 &  \\
\bottomrule
\end{tabular}

\section{基础题}

\subsection{题目与链接}

\begin{itemize}
\item 选择题目：
\item 题目链接：
\end{itemize}

\subsection{机器环境}

\begin{tabular}{ll}
\toprule
项目 & 内容 \\
\midrule
机器来源 & 本地电脑 / WSL / 虚拟机 / 自备服务器 / 其他 \\
操作系统 &  \\
CPU &  \\
内存 &  \\
GPU &  \\
编译器 &  \\
MPI &  \\
BLAS / 数学库 &  \\
关键依赖版本 &  \\
\bottomrule
\end{tabular}

\subsection{Baseline 与测试配置}

\begin{tabular}{llllll}
\toprule
配置编号 & 运行命令 & 关键参数 & 运行时间 & 性能指标 & 正确性结果 \\
\midrule
1 &  &  &  &  &  \\
2 &  &  &  &  &  \\
3 &  &  &  &  &  \\
\bottomrule
\end{tabular}

\subsection{基础题结论}

说明哪组参数或配置表现最好，以及判断依据。

\section{大题}

\subsection{题目理解}

\begin{itemize}
\item 选择题目：
\item 题目链接：
\item 大题主要内容：
\item 性能目标：
\item 正确性判断方式：
\end{itemize}

\subsection{机器环境}

\begin{tabular}{ll}
\toprule
项目 & 内容 \\
\midrule
机器来源 & 本地电脑 / WSL / 虚拟机 / 自备服务器 / Colab / 其他 \\
操作系统 &  \\
CPU &  \\
内存 &  \\
GPU &  \\
编译器 &  \\
MPI &  \\
Python &  \\
CUDA / ROCm &  \\
关键依赖版本 &  \\
\bottomrule
\end{tabular}

\subsection{Baseline}

\begin{lstlisting}[language=bash]

\end{lstlisting}

\begin{tabular}{ll}
\toprule
项目 & 内容 \\
\midrule
运行时间 &  \\
性能指标 &  \\
正确性结果 &  \\
输出文件 &  \\
日志文件 &  \\
\bottomrule
\end{tabular}

\subsection{性能分析}

说明使用了哪些工具或方法进行分析，并给出至少一个有证据支持的结论。

\subsection{优化或工程改进}

\begin{tabular}{llll}
\toprule
编号 & 修改内容 & 修改原因 & 预期效果 \\
\midrule
1 &  &  &  \\
2 &  &  &  \\
3 &  &  &  \\
\bottomrule
\end{tabular}

\subsection{优化后结果}

\begin{tabular}{lllll}
\toprule
实验 & 运行时间 & 性能指标 & 正确性结果 & 说明 \\
\midrule
Baseline &  &  &  &  \\
优化后 &  &  &  &  \\
\bottomrule
\end{tabular}

\subsection{正确性验证}

说明优化后结果如何验证，是否与 Baseline 或官方要求一致。

\section{遇到的问题和解决方法}

\begin{tabular}{lll}
\toprule
问题 & 现象 & 解决方法 \\
\midrule
 &  &  \\
\bottomrule
\end{tabular}

\section{总结}

总结基础题和大题的完成情况、有效优化、无效优化、对题目的理解，以及后续还可以继续优化的方向。

\end{document}
